\clearpage
\fancyhead[L]{\small ASTRONOMISCHES JAHRBUCH.}
\hypertarget{native-610-lower}{}
\begin{center}\small Göttingische gelehrte Anzeigen. Stück 81 u. 82. Seite 801..810. 1818 Mai 21.\\[3pt]\rule{.72\linewidth}{.3pt}\end{center}

\emph{Astronomisches Jahrbuch für das Jahr {\upshape 1820}, nebst einer Sammlung der neuesten in die astronomischen Wissenschaften einschlagenden Abhandlungen, Beobachtungen und Nachrichten. Mit Genehmhaltung der k. Acad. d. W. berechnet und herausgegeben von} Dr. I. E. \textsc{Bode}, königl. Astronom u. s. w. Berlin. Bei dem Verfasser, und in Comm. bei F. Dümmler. 256 Seiten in Octav, nebst einer Kupfertafel.

Das Jahr 1820 zeichnet sich durch mehrere merkwürdige astronomische Erscheinungen aus. Ausser der grossen Sonnenfinsterniss vom 7. September, welche auch in unserer Gegend ringförmig sein wird, kommen auch mehrere Planetenbedeckungen vom Monde vor, unter denen eine des Mars am 28. Januar, und zwei des Jupiter am 4. Juni und 18. October in Berlin sichtbar sein werden. --- Den Anfang der Zusätze machen die astronomischen Beobachtungen auf der Berliner Sternwarte im Jahre 1816, wovon wir die Beobachtung der Sonnenfinsterniss vom 19. November, und die der Mondfinsterniss vom 4. December auszeichnen. --- Beobachtungen in Halberstadt vom Hauptmann von \textsc{Wahl}. Die Breite von Halberstadt (Moritzplan), wurde durch Sonnenbeobachtungen mit einem 18zolligen \textsc{Troughton}'schen Spiegelkreise gefunden $51^\circ\,54'\,5''9$. Die Länge berechnet Hr. \textsc{von Wahl} aus einer Sonnenfinsterniss und einer Bedeckung des Aldebaran $1'\,17''6$ in Zeit östlich von der Seeberger Sternwarte; die hier nach Herrn \textsc{von Wahl}'s Rechnung bestimmten Conjunctionszeiten für mehrere verglichene Oerter weichen einige Secunden von den Resultaten anderer Berechner ab. --- Sichtbare Lichtveränderungen Algols in den Jahren 1817---1819 berechnet vom Prof. \textsc{Wurm}. --- Astronomische Anzeigen und Beobachtungen des Gegenscheins des Jupiter 1816 vom Astronomen \textsc{Derflinger} zu Kremsmünster. --- Astronomische Bemerkungen und Berechnung des Gegenscheins der Vesta 1815 und der Sonnenfinsterniss von 1820, vom Doctor \textsc{Gerling} in Cassel. Dieser Artikel enthält auch eine Berichtigung eines im Jahrbuch für 1819 vorkommenden übereilten Urtheils über die atmosphärische Refraction. --- Ueber das \textsc{Kepler}'sche Problem vom Staatsrath und Ritter \textsc{von Schubert}. Dieser Aufsatz ist gewissermassen eine ausführliche Entwicklung der \textsc{Laplace}'schen Auflösung der Aufgabe, durch eine nach den Sinussen der Vielfachen der mittlern Anomalie fortschreitende unendliche Reihe, deren Coëfficienten selbst wieder in der Form unendlicher nach den Potenzen der Excentricität fortlaufender Reihen erscheinen; in einem weiterhin folgenden Nachtrage hat der Verf. die allerdings sehr mühsame Entwicklung bis zur 13ten Potenz der Excentricität getrieben. In so fern diese Entwicklung in mathematischer Rücksicht interessant ist, muss man dem Verf. dafür Dank wissen, wenn man gleich über den astronomischen Werth der Auflösung durch eine solche unendliche Reihe, verglichen mit sogenannten indirecten Auflösungen, ganz anders urtheilt als jener, und gerade dasjenige, was er zur Empfehlung der erstern auf Kosten der andern geltend \SourceLeaf{611}machen will, nur im umgekehrten Sinn anwendbar findet. Denn gerade bei dem Gebrauch der Reihe kann man die Genauigkeit nicht weiter treiben, als man die Werthe der numerischen Coëfficienten entwickelt hat, während die Genauigkeit der sogenannten indirecten Methode lediglich durch die Genauigkeit der gebrauchten Sinustafeln bedingt wird. Dass letztere, wenn sie zweckmässig eingerichtet sind, bei etwas beträchtlichen Excentricitäten, ohne Vergleich bequemer sind, ist ohnehin bekannt genug. --- Astronomische Beobachtungen auf der Sternwarte zu Wien 1816 von \textsc{Triesnecker} und \textsc{Bürg}; Beobachtungen von Jupiterstrabanten-Verfinsterungen, Sternbedeckungen und einigen Planetenoppositionen, der Mondfinsterniss vom 9. Juni und der Sonnenfinsterniss vom 18. November. --- Ueber die Verbesserung des Mittagsfernrohrs vom Prof. \textsc{Littrow}. Dieser Aufsatz beschäftigt sich mit der Aufgabe, die Lage des Mittagsfernrohrs aus beobachteten Durchgängen zu bestimmen, wenn die Fehler eine endliche Grösse haben. Eigentlich hat diese Aufgabe mehr ein mathematisches, als ein astronomisches Interesse, da in der Ausübung nur solche Fehler vorkommen, die als unendlich klein zu betrachten sind. Ueberdies ist die Anwendung dieser Methode von der Kenntniss der wahren Culminationszeiten abhängig, die man nur mit Hülfe eines andern Instruments erhalten kann; zu einer selbstständigen Berichtigung des Mittagsfernrohrs dürfte man nur die Unterschiede der wahren Culminationszeiten als bekannt ansehen, und muss dann die Axe des Instruments als berichtigt, oder die Neigung als gegeben voraussetzen. Freilich ist auch dies Verfahren nur für untergeordnete Mittagsfernröhre passend, und zur Berichtigung eines Instruments vom ersten Range eine von der Kenntniss der Rectascensionen abhängige Berichtigung unzulässig. --- Aus einem Schreiben desselben Astronomen wird ein Auszug einer zur Bestimmung der die Bewegung der Erde störenden Planetenmassen abzweckenden Untersuchung gegeben, über welche wir, da sie in der Zeitschrift für Astronomie ausführlich abgedruckt ist, uns einige Bemerkungen erlauben. Die Vergleichung mit letzterer zeigt, dass in vorliegendem Auszuge ein Druckfehler die Unsicherheit (richtiger den wahrscheinlichen Fehler) der Venusmasse zehnfach zu gross gemacht hat. Der Verf. hat 189 Greenwicher Sonnenbeobachtungen mit den \textsc{von Zach}'schen Sonnentafeln verglichen und durch gruppenweise Zusammenfassung, 43 Bedingungsgleichungen entwickelt, die 11 unbekannte Grössen enthalten Er schränkt sich jedoch darauf ein, nur von einer derselben, nemlich von den Verbesserungen der Epoche, der Venusmasse, der Marsmasse und der Mondsmasse die Werthe nach der Methode der kleinsten Quadrate zu bestimmen. Wir sehen jedoch den Grund nicht recht ein, warum eben die Epoche von den elliptischen Elementen allein, vorzugsweise vor der Länge des Perigäum und der Excentricität berücksichtigt ist. Auch bedürfte wohl der Ausdruck, dass die Epoche nicht leicht mit einiger Schärfe anzugeben sei, einer Berichtigung. Ferner ist das Urtheil (Zeitschrift S. 285), dass die Nutation aus Sonnenbeobachtungen mit mehr Sicherheit als die Venusmasse zu bestimmen sei, unhaltbar, wenigstens wenn von Sonnenlängen die Rede ist, die durch beobachtete Rectascensionen bestimmt sind, denn das unmittelbar Beobachtete besteht aus Rectascensionsunterschieden der Sonne und der Fundamentalfixsterne, welche Unterschiede nur schwach von der Nutation afficirt werden; die Sonnenlängen selbst involviren schon schon die bei den Fixsternen berechnete Nutation, und man wird daher immer nur fast genau dieselbe Nutation wieder finden müssen, die man bei der Rechnung zum Grunde gelegt hatte. Was nun endlich die Resultate selbst betrifft, die Hr. \textsc{Littrow} für die erwähnten vier Grössen selbst herausgebracht hat, so ist eine Wiederholung der Rechnung deswegen sehr zu wünschen, weil in derselben mehrere Fehler sichtbar sind, von denen sich nicht wohl entscheiden lässt, ob es Druckfehler, Schreibfehler oder Rechnungsfehler sind. Die Natur der Methode der kleinsten Quadrate bringt es mit sich, dass die Coëfficienten von $\mu'''$ in der ersten und von $\mu$ in der dritten Gleichung S. 292 gleich sein müssen, da sie hier ganz verschieden sind; auch die Coëfficienten von $\mu'''$ in der zweiten und von $\mu'$ in der dritten Gleichung sollen ganz dieselben sein. Die Summen der Quadrate der Fehler, vor Anbringung \SourceLeaf{612}der gefundenen Verbesserung fand Hr. \textsc{Littrow} 973, nach derselben 727; es scheint uns aber, dass diese Verminderung beträchtlich stärker hätte ausfallen müssen. Dass es zweckmässiger sei, die Fehler mit ihren Zeichen selbst, als deren Quadrate zu addiren, ist nicht anzunehmen; denn dass jene Summe verschwindet, beweist \emph{blos} für die Wirksamkeit der Verbesserung der Epoche, und gar nichts für die andern Verbesserungen. Wir wünschen, dass der Verf. diese Bemerkungen, als ein Zeichen der Aufmerksamkeit, womit wir seiner schätzbaren Untersuchung gefolgt sind, ansehen, und diese, seinem Versprechen zufolge, bald nach einem grössern Maassstabe ausführen möge. --- Astronomische Beobachtungen vom Jahr 1816 auf der Prager und auf der Wilnaer Sternwarte; unter jenen befindet sich die vollständige Beobachtung der Sonnenfinsterniss vom 18. November, und die Opposition der Ceres, unter diesen die Opposition der Pallas. --- Vorzüglich schätzbar sind die Beobachtungen vom Prof. \textsc{Bessel} auf der Königsberger Sternwarte, die sechs Planetenoppositionen, fünf Sternbedeckungen, die Mondfinsterniss vom 4. Dec. 1816, die Zenithdistanzen der Sonne in beiden Solstitien, und die gerade Aufsteigung des Polarsterns befassen. --- Bemerkungen bei Gelegenheit der grossen Sonnenfinsterniss am 19. Nov. 1816. Da eine grosse Sonnenfinsterniss eine von den seltenen Gelegenheiten ist, wo der Zustand der Witterung von einer grossen Anzahl von Oertern zur öffentlichen Kenntniss kommt, so hatte der Herausgeber den glücklichen Gedanken, alle die Oerter zusammenzustellen, wo ein heiterer Himmel die vollständige Beobachtung erlaubte, ferner die, wo abwechselndes Wetter nur unterbrochen etwas von der Finsterniss zu bemerken möglich machte, und endlich die, wo ganz trüber Himmel gar nichts davon zu Gesicht kommen liess. Der Erfolg zeigt, dass alle diese Oerter ohne Zusammenhang und Regelmässigkeit bunt durch einander liegen, und ist daher sehr geeignet, den Glauben an specielle Wetterprophezeiungen etwas niederzuschlagen. Noch fügt der Verf. einige Bemerkungen über die Erfahrung bei, dass gewöhnlich bei totalen Finsternissen die Dunkelheit nicht so gross ist, wie man erwartet; was er aber als Versuch einer Erklärung des um den Mond bei totalen Finsternissen zuweilen bemerkten Ringes anführt, ist uns nicht ganz klar geworden. --- Beobachtungen der Jupiterstrabantenverfinsterungen und Sternbedeckungen auf der Greenwicher Sternwarte von 1811---1814, und die eben daselbst angestellten Beobachtungen des grossen Cometen von 1811. --- Neue Elemente der Vestabahn vom Prof. \textsc{Gerling}. --- Astronomische Beobachtungen auf der Göttinger Sternwarte vom Prof. \textsc{Gauss}; Mondfinsterniss vom 4. Decemb. 1816, Polhöhe der neuen Sternwarte aus 158 Beobachtungen des Nordsterns, und Beobachtungen des Sommersolstitium von 1817. Bei diesen Beobachtungen ergab sich dasselbe Resultat, welches seit mehreren Jahren die Astronomen in Verlegenheit gesetzt hat, dass nemlich die Polhöhe aus Circumpolarsternen um mehrere Secunden grösser ausfällt, als aus Sonnenbeobachtungen, obgleich die erwähnten Beobachtungen ohne das am Objectiv sonst angesteckte und von einigen als mögliche Ursache jenes Unterschiedes in Verdacht gezogene Gegengewicht angestellt waren. Seitdem diese Beobachtungen, welche die Unzulässigkeit dieser Vermuthung beweisen, durch die Zeitschrift für Astronomie bekannt gemacht sind, hat Hr. Prof. \textsc{Bohnenberger} eine neue Hypothese zur Erklärung des räthselhaften Phänomens aufgestellt, auf deren Veranlassung an dem \textsc{Reichenbach}'schen Kreise, mit welchem obige Beobachtungen angestellt sind, eine Abänderung angebracht ist, deren Wirkung bei den Beobachtungen des letzten Wintersolstitium der Richtigkeit der \textsc{Bohnenberger}'schen Hypothese günstig zu sein scheint; ob aber diese hinlänglich ist, das Phänomen \emph{ganz} zu erklären, oder ob nicht vielmehr eine Conspiration \emph{mehrerer} Ursachen angenommen werden müsse, wird an einem andern Ort in Untersuchung gezogen werden. --- Beobachtung der Sonnenfinsterniss vom 19. Nov. 1816 zu Glatz vom General \textsc{von Lindener}. --- Nachrichten über des verstorbenen \textsc{Triesnecker}'s Lebensumstände vom Prof. \textsc{Bürg}. --- Bestimmungen für den Polarstern vom Baron \textsc{von Lindenau}. Die sehr verdienstliche Discussion von achthundert Beobachtungen des Polarsterns hat mehrere interessante Resultate gegeben; eine neue Bestimmung der \SourceLeaf{613}Aberrationsconstante zu $20''44861$ (wahrscheinlicher Fehler nach der \textsc{Gauss-Laplace}'schen Wahrscheinlichkeitstheorie $\pm0''032$), der Nutationsconstante zu $8''97707$ (wahrscheinlicher Fehler $\pm0''04421$), und der Parallaxe des Polarsterns zu $0''14444$ (wahrscheinlicher Fehler $\pm0''05568$). Aus der Nutationsconstante verbunden mit der Präcessionsconstante hat Hr. \textsc{von Lindenau} vermittelst einiger Formeln der Mécanique Céleste und einer ihm von \textsc{Gauss} mitgetheilten Bedingungsgleichung, auch noch die Abplattung der Erde und die Mondsmasse abzuleiten versucht, und für erstere nahe $\tfrac{1}{816}$ gefunden; allein es haben sich in die angewandten Formeln und in die Rechnung mehrere Fehler eingeschlichen, deren Berichtigung durch Hrn \textsc{von Lindenau} selbst an einem schicklichen Orte zu erwarten ist. Es ist ein blosser Zufall, dass obiges Resultat der Wahrheit so nahe kommt; wenn alle Fehler gehörig verbessert werden, ergibt sich die Abplattung etwa zu $\tfrac{1}{203}$, wobei jedoch bemerkt werden muss, dass die von \textsc{Gauss} gegebene Bedingungsgleichung (aus der, so wie sie hier abgedruckt ist, der Coëfficient $\tfrac{3}{5}$ weggelassen werden muss) in so fern hypothetisch ist, als dabei die Aehnlichkeit der Lagen von gleicher Dichtigkeit im Erdsphäroid vorausgesetzt ist. --- Beobachtete Oppositionen vom Jupiter und Uranus 1817 auf der Seeberger Sternwarte vom Adjunct \textsc{Enke}; letztere ist besonders deswegen wichtig, weil Uranus noch nie so nahe bei seinem Knoten beobachtet worden ist. --- Ueber die Verbesserung einer schon beiläufig bekannten Cometenbahn, vom Dr. \textsc{Olbers}. Es wird hier ein unpassendes Urtheil, welches ein Astronom über die von \textsc{Olbers} zu dem angegebenen Zweck vorgeschlagene Methode gefällt hatte, berichtigt, und zugleich gezeigt, wie diese Methode durch Hinzufügung einer vierten Hypothese, zur Bestimmung der elliptischen Bahn eingerichtet werden könne, ein brauchbares Verfahren, welches neben andern empfohlen zu werden verdient. Der Verf. bemerkt übrigens mit Recht, dass es eine nicht zu billigende Verschwendung von Kraft und Zeit sein würde, wenn man bei \emph{allen} Cometen die Bahn elliptisch zu berechnen unternehmen wollte. Dies Urtheil darf jedoch niemanden, der die Kräfte dazu hat, abschrecken, die Bahn eines Cometen, dessen beobachtete Bewegung entschieden von der Parabel abweicht, mit aller Sorgfalt zu bestimmen; auch wird der hochverdiente Verf. es gewiss nicht missbilligen, wenn angehende Astronomen, ihre Kräfte zu einem Geschäfte, welches noch nicht Jedermanns Sache ist, zu üben, sich etwa auch an einem Cometen versuchen, bei dem der Erfolg am Ende nur eine in ziemlich schwankende Grenzen eingeschlossene Abweichung von der Parabel zeigt. --- Neue Berechnung der Säcularänderungen der Elemente der Erdbahn, vom Prof. \textsc{Nicolai}. Diese sehr verdienstliche Arbeit wurde durch die Vermuthung veranlasst, dass der Unterschied zwischen der Bestimmung der Venusmasse, welche Hr. \textsc{von Lindenau} durch die Mercurstheorie gefunden hat, und derjenigen, welche aus der Abnahme der Schiefe der Ekliptik nach \textsc{Laplace}'s Theorie gefolgert ist, in dem Umstande seinen Grund haben könnte, dass letzter blos eine genäherte Bestimmung der Säcularänderungen gibt. Der Verf. hat daher die Säcularänderungen sämmtlicher Elemente der Erdbahn nach der neuen (noch nicht öffentlich bekannt gemachten) \textsc{Gauss}'schen Methode, die vollkommen streng ist, mit grösster Sorgfalt neu berechnet; der Erfolg hat indessen jene Vermuthung nicht bestätigt, da die neuen Resultate von den ältern doch nicht bedeutend verschieden ausgefallen sind. --- Von demselben Astronomen Beobachtung der Oppositionen der Vesta, Pallas, des Jupiter und des Uranus 1817 auf der Mannheimer Sternwarte. --- Geometrischer Lauf 1817---1818 der Vesta berechnet vom Prof. \textsc{Gerling} in Marburg, und der Pallas vom Dr. \textsc{Tittel} in Göttingen. --- Auszug aus einem Schreiben des Dr. \textsc{Struve} in Dorpat gibt unter andern Nachricht von einer vorzunehmenden trigonometrischen Vermessung Lieflands. --- Preisverzeichniss astronomischer Instrumente im Institut des Hrn \textsc{von Utzschneider} in Benediktbeuern. --- Die vermischten astronomischen Beobachtungen, Bemerkungen und Nachrichten, welche diesen Jahrgang beschliessen, enthalten auch noch mehrere astronomische und geographische Bestimmungen, die wir des beschränkten Raumes wegen hier nicht einzeln namhaft machen können.
